\subsection{Sharpness audit via a state-constrained control problem}
\label{subsec:approx-flat-sharpness}

The preceding estimates are upper bounds, but their logical status is not
the same in the two signs of $\ellzero$. We now formulate the trace problem
as an exact one-dimensional control system. This formulation proves
sharpness for $\ellzero\leqslant0$, detects a strict gap in the positive-
$\ellzero$ estimate of Corollary~\ref{cor:optimized-approx-flat-span}, and
identifies the residual optimization problem without confusing a formal
trace with a solution of the two-dimensional field equation.

For $\ellzero>0$, define
\begin{equation}
\mathfrak g(u,y,s):=
\frac{2(1+s)\bigl[(1-s)-y(1-u^2s)\bigr]}
{2+y(1-u^2)}.
\label{eq:positive-l-control-drift}
\end{equation}
The exact equatorial identity in
Lemma~\ref{lem:exact-equatorial-log-identity} is equivalent to
\begin{equation}
\dot u=su,\qquad
\dot y=(1-s)y,\qquad
\dot s=\mathfrak g(u,y,s)+w,
\qquad w:=-Z\geqslant0,
\label{eq:positive-l-control-system}
\end{equation}
where a dot denotes $d/dt$. The state constraints are
\begin{equation}
0<u<1,\qquad
0<y\leqslant\frac{2}{1+u^2},\qquad
|s|\leqslant\varepsilon.
\label{eq:positive-l-control-state}
\end{equation}
Thus the logarithmic span is precisely the exit time of a
state-constrained system with the one-sided transverse-curvature control
$w\geqslant0$. In particular, the control can increase $s$ relative to the
zero-curvature drift, but it cannot decrease it.

There is also a geometric version of the same reduction. Set
$X=r^2$ and let
$f(X)=\Phi_{\ellzero}(r,ru(r))$ be the auxiliary trace. A direct
calculation gives
\begin{equation}
\frac{d^2f}{dX^2}
=\frac{uA}{4r^3}\bigl(\dot s-\mathfrak g(u,y,s)\bigr)
=\frac{uA}{4r^3}w,
\qquad A=2+y(1-u^2)>0.
\label{eq:control-convexity-equivalence}
\end{equation}
Consequently, vertical localization is exactly convexity of the auxiliary
equatorial datum as a function of $r^2$. This is a useful structural
interpretation of the control; it is not an additional inequality.

\begin{lemma}[Analytic lifting of strict trace controls]
\label{lem:analytic-lifting-trace-controls}
Let $u$ be positive and real analytic near a compact radial interval away
from the axis. Suppose that the inequalities in
Eq.~\eqref{eq:positive-l-control-state} are strict, that $A>0$, and that
$\dot s>\mathfrak g(u,y,s)$. Then the trace lifts to a reflection-symmetric
real-analytic solution of the exact constant-$\ell$ VFE in a neighbourhood
of the interval, with $v(r,0)=u(r)$, $v_z(r,0)=0$, positive density, and
$v_{zz}(r,0)<0$.
\end{lemma}

\begin{proof}
Prescribe analytic Cauchy data for the linear auxiliary equation by
$$
\mathcal F(r,0)=\Phi_{\ellzero}(r,ru(r)),
\qquad \mathcal F_z(r,0)=0.
$$
The Cauchy--Kowalevski theorem applied to
$\mathcal F_{zz}=-\mathcal F_{rr}+r^{-1}\mathcal F_r$ gives a unique
analytic solution near each point of the interval. Uniqueness makes the
local solutions agree on overlaps, and compactness gives a common collar.
Since $A=\partial_\eta\Phi_{\ellzero}>0$, the analytic implicit-function
theorem recovers a unique analytic field $\eta$ on the same inverse sheet.
The chain-rule identity underlying Eq.~\eqref{eq:GS-operator} then gives
the exact VFE. On the trace,
$Z=\mathfrak g-\dot s<0$, so $v_{zz}=uZ/r^2<0$. Strict density and
subluminality persist in a smaller collar by continuity.
\end{proof}

The lemma shows that strict analytic controls are genuine local PDE
solutions. It does not make the elliptic Cauchy problem stable, nor does it
provide an axis, a vacuum exterior, or a global completion.

We next improve the positive-$\ellzero$ comparison estimate. The result is
deliberately explicit because it proves that the previously optimized
bound is not the true optimum.

\begin{theorem}[Strict positive-$\ell$ gap and a refined span bound]
\label{thm:refined-positive-l-span}
Assume the hypotheses of
Theorem~\ref{thm:quantitative-approx-flat-no-go}, let
$0<\varepsilon<1$, suppose $\ellzero>0$, and write
$v_{\mathrm{min}}:=\underline v\in(0,1)$. For
$\alpha\in(0,1-\varepsilon)$, define
\begin{align}
c_+(\varepsilon,\alpha;v_{\mathrm{min}})
&:=\frac{2(1+\varepsilon)
\bigl[1-\varepsilon-\alpha(1-v_{\mathrm{min}}^2\varepsilon)\bigr]}
{2+\alpha(1-v_{\mathrm{min}}^2)},
\label{eq:c-plus-refined}\\
c_-(\varepsilon,\alpha;v_{\mathrm{min}})
&:=
\begin{cases}
\displaystyle
\frac{2(1-\varepsilon)
\bigl[1+\varepsilon-\alpha(1+v_{\mathrm{min}}^2\varepsilon)\bigr]}
{2+\alpha(1-v_{\mathrm{min}}^2)},
&\displaystyle
\alpha\leqslant\frac{1-\varepsilon}{1+\varepsilon},\\[3mm]
(1-\varepsilon^2)(1-\alpha),
&\displaystyle
\alpha>\frac{1-\varepsilon}{1+\varepsilon},
\end{cases}
\label{eq:c-minus-refined}\\
c_*(\varepsilon,\alpha;v_{\mathrm{min}})&:=\min\{c_+,c_-\}.
\label{eq:c-star-refined}
\end{align}
Then $c_*>0$ and
\begin{equation}
\Lambda\leqslant
\mathcal B_*(\varepsilon,\alpha;v_{\mathrm{min}}):=
\frac{2\varepsilon}{c_*(\varepsilon,\alpha;v_{\mathrm{min}})}
+\frac{1}{1-\varepsilon}
\log\!\left[\frac{2}{(1+v_{\mathrm{min}}^2)\alpha}\right].
\label{eq:refined-positive-l-bound}
\end{equation}
Moreover, if
\begin{equation}
c_0(\varepsilon,\alpha):=
\frac{2(1+\varepsilon)(1-\varepsilon-\alpha)}{2+\alpha},
\label{eq:c-zero-old-comparison}
\end{equation}
then
\begin{equation}
c_*(\varepsilon,\alpha;v_{\mathrm{min}})>c_0(\varepsilon,\alpha),
\qquad
\mathcal B_*(\varepsilon,\alpha;v_{\mathrm{min}})
<\mathcal B_+(\varepsilon,\alpha;v_{\mathrm{min}}).
\label{eq:strict-positive-l-gap}
\end{equation}
In particular, at the optimizer $\alpha_\varepsilon$ in
Eq.~\eqref{eq:alpha-optimal},
\begin{equation}
\begin{aligned}
\Lambda
&\leqslant\inf_{0<\alpha<1-\varepsilon}
\mathcal B_+(\varepsilon,\alpha;v_{\mathrm{min}})
\\
&\quad-\delta(\varepsilon,v_{\mathrm{min}}),\\
\delta(\varepsilon,v_{\mathrm{min}})
&:=2\varepsilon\left[
\frac1{c_0(\varepsilon,\alpha_\varepsilon)}
\right.\\
&\left.\hspace{2.4cm}
-\frac1{c_*(\varepsilon,\alpha_\varepsilon;v_{\mathrm{min}})}
\right]>0.
\end{aligned}
\label{eq:strict-gap-from-old-optimum}
\end{equation}
Thus the earlier positive-$\ell$ estimate is separated from the true
trace optimum for every fixed $\varepsilon>0$ and
$v_{\mathrm{min}}>0$; it is neither
attained nor asymptotically attained with those parameters fixed.
\end{theorem}

\begin{proof}
On the portion of the trace where $y\leqslant\alpha$, the master
inequality reads $\dot s\geqslant\mathfrak g(u,y,s)$. Put $x=u^2$.
Differentiation of the drift gives
\begin{align}
\partial_y\mathfrak g
&=\frac{2(1+s)(sx+s+x-3)}{[2+y(1-x)]^2},
\label{eq:drift-y-derivative}\\
\partial_x\mathfrak g
&=\frac{2y(1+s)[(1+s)-y(1-s)]}{[2+y(1-x)]^2}.
\label{eq:drift-x-derivative}
\end{align}
For $|s|\leqslant\varepsilon<1$ and $0\leqslant x<1$, the numerator in
Eq.~\eqref{eq:drift-y-derivative} is negative. Hence the minimum over
$0<y\leqslant\alpha$ occurs at $y=\alpha$. For fixed $x$, the derivative
with respect to $s$ has the sign of
$$
2s(x\alpha-1)+\alpha(x-1),
$$
which is strictly decreasing in $s$. Any interior stationary point is
therefore a maximum, and the minimum over
$[-\varepsilon,\varepsilon]$ occurs at an endpoint. At
$s=\varepsilon$, Eq.~\eqref{eq:drift-x-derivative} is positive, so the
minimum over $v_{\mathrm{min}}^2\leqslant x<1$ is $c_+$ at
$x=v_{\mathrm{min}}^2$. At
$s=-\varepsilon$, its sign is that of
$(1-\varepsilon)-\alpha(1+\varepsilon)$. This gives the two cases in
Eq.~\eqref{eq:c-minus-refined}; in the second case the infimum is the
limit $x\uparrow1$. It follows that
$$
\dot s\geqslant c_*(\varepsilon,\alpha;v_{\mathrm{min}})>0
$$
whenever $y\leqslant\alpha$. Since $s$ remains in an interval of width
$2\varepsilon$, this initial portion has length at most
$2\varepsilon/c_*$. On the remaining portion,
$d\log y/dt=1-s\geqslant1-\varepsilon$, while density positivity gives
$y\leqslant2/(1+u^2)\leqslant2/(1+v_{\mathrm{min}}^2)$. Its length is
therefore at most
the logarithmic term in Eq.~\eqref{eq:refined-positive-l-bound}.

It remains to prove that the improvement is strict. Direct subtraction
gives
\begin{equation}
c_+-c_0=
\frac{2\alpha v_{\mathrm{min}}^2(1+\varepsilon)
\bigl[1+\varepsilon-\alpha(1-\varepsilon)\bigr]}
{(2+\alpha)[2+\alpha(1-v_{\mathrm{min}}^2)]}>0.
\label{eq:c-plus-strict-difference}
\end{equation}
If
$\alpha\leqslant(1-\varepsilon)/(1+\varepsilon)$, monotonicity in $x$
at $s=-\varepsilon$ gives
$c_-\geqslant\mathfrak g(0,\alpha,-\varepsilon)$, and
\begin{equation}
\mathfrak g(0,\alpha,-\varepsilon)-c_0
=\frac{4\alpha\varepsilon}{2+\alpha}>0.
\label{eq:c-minus-first-strict-difference}
\end{equation}
In the other case,
\begin{equation}
c_--c_0=
\frac{\alpha(1+\varepsilon)
\bigl[1+\varepsilon-\alpha(1-\varepsilon)\bigr]}
{2+\alpha}>0.
\label{eq:c-minus-second-strict-difference}
\end{equation}
Thus $c_*>c_0$. The first term in the old bound is
$2\varepsilon/c_0$, whereas its logarithmic term is unchanged, proving
Eqs.~\eqref{eq:strict-positive-l-gap} and
\eqref{eq:strict-gap-from-old-optimum}.
\end{proof}

The refined comparison is still not asserted to be optimal. It minimizes
the drift over a rectangular state enclosure and then adds two exit-time
estimates. The exact positive-$\ellzero$, positive-$\varepsilon$ supremum
requires the coupled evolution of $(u,y,s)$, including possible arcs on
$s=\varepsilon$, and remains open. Equation~\eqref{eq:positive-l-control-system}
is the precise problem to be solved; any proposed extremal must also obey
the density boundary and the tangency condition
$\mathfrak g(u,y,\varepsilon)\leqslant0$ on a sustained upper-slope arc.

The limiting exactly flat problem is more rigid and can be settled.

\begin{proposition}[Sharp exact-flat limit for positive $\ell$]
\label{prop:sharp-positive-l-exact-flat-limit}
Fix $0<v_0<1$ and $\ellzero>0$. Among real-analytic exactly flat traces
with strict positive density and $v_{zz}<0$, the supremum of the radial
ratio is
\begin{equation}
\sup\frac{R_2}{R_1}=\frac{2}{1+v_0^2}.
\label{eq:positive-l-exact-flat-supremum}
\end{equation}
The supremum is approached by genuine reflection-symmetric analytic
solutions of the two-dimensional VFE. It is not attained within the
strict-margin class.
\end{proposition}

\begin{proof}
The upper bound is Eq.~\eqref{eq:flat-trace-span-ratio}. Conversely, choose
closed intervals compactly contained in
$$
\frac{v_0}{\ellzero}<r<
\frac{2v_0}{\ellzero(1+v_0^2)}
$$
whose endpoint ratio tends to $2/(1+v_0^2)$. The analytic trace
$u\equiv v_0$ has strict density because
$\ellzero r<2v_0/(1+v_0^2)$, while
Eq.~\eqref{eq:vzz-x-form} gives $v_{zz}<0$ because
$\ellzero r>v_0$. Lemma~\ref{lem:analytic-lifting-trace-controls}, or
equivalently the analytic auxiliary Cauchy construction used in
Theorem~\ref{thm:local-flat-trace-existence}, supplies a genuine analytic
solution near each interval. The two strict inequalities prevent an
endpoint from attaining the limiting ratio.
\end{proof}

Finally, the nonpositive-$\ellzero$ result admits a stronger sharpness
statement than attainment by the endpoint $\ellzero=0$ alone.

\begin{proposition}[Sharpness throughout the nonpositive-$\ell$ closure]
\label{prop:strict-negative-l-asymptotic-sharpness}
For every $0<\varepsilon<1$, the bound
$\Lambda\leqslant2\operatorname{artanh}\varepsilon$ is attained at
$\ellzero=0$ by Eq.~\eqref{eq:sharp-span-extremal}. If
$\ellzero<0$, every nondegenerate trace satisfying the same hypotheses
obeys the strict inequality
\begin{equation}
\Lambda<2\operatorname{artanh}\varepsilon.
\label{eq:strict-negative-l-span}
\end{equation}
Nevertheless, the right-hand side is the supremum over the strict class
$\ellzero<0$: there is a sequence of genuine analytic, reflection-symmetric,
vertically neutral solutions with $\ellzero\uparrow0$ whose spans converge
to $2\operatorname{artanh}\varepsilon$.
\end{proposition}

\begin{proof}
For $\ellzero=0$, let
$u(t)=a\cosh(t-t_0)$ and $s(t)=\tanh(t-t_0)$. Then
$\dot s=1-s^2$, and
$$
\eta=ru=\frac a2\left(\frac{r^2}{r_0}+r_0\right)
$$
satisfies $\mathcal L(2\eta)=0$. Hence the solution is independent of
$z$, has $v_{zz}=0$, and traverses from $s=-\varepsilon$ to
$s=\varepsilon$ in exactly
$2\operatorname{artanh}\varepsilon$. Taking $a$ sufficiently small gives
strict subluminality and positive density on the compact annulus.

For $\ellzero<0$, the exact factorization
\eqref{eq:sharp-negative-l-factorization} is strictly positive because
$y<0$, $1+s>0$, and
$3-u^2-(1+u^2)s>0$. Since $A>0$, it follows that
$\dot s>1-s^2$ pointwise. Integration proves
Eq.~\eqref{eq:strict-negative-l-span}.

To prove asymptotic sharpness, keep the preceding radial auxiliary datum
$\mathcal F_0=2\eta_0=c_2r^2+c_0$ fixed. For all sufficiently small
negative $\ellzero$, the analytic implicit-function theorem solves
$$
\Phi_{\ellzero}(r,\eta_{\ellzero}(r))=\mathcal F_0(r)
$$
on a common annulus, because
$\partial_\eta\Phi_0=2$. The resulting $\eta_{\ellzero}$ is independent
of $z$ and solves the exact VFE. Its speed, slope, and density converge
uniformly with their required derivatives to the strict-margin
$\ellzero=0$ solution. The two transverse crossings
$s=-\varepsilon$ and $s=\varepsilon$ persist by the implicit-function
theorem because $\dot s=1-\varepsilon^2>0$ at the limiting crossings.
Their time separation therefore converges to
$2\operatorname{artanh}\varepsilon$. This construction proves sharpness
within genuine analytic PDE solutions, not merely within formal trace
controls.
\end{proof}
